\chapter{Conclusion}
\quad This thesis has investigated a MARL framework for FLSO to address GW load balancing in LEO NTN. Focusing on regenerative payload architectures where backhaul mobility is decoupled from user access, we formulated the FLSO decision problem as a decentralized MMDP. The proposed framework, built upon realistic assumptions including a Walker-delta constellation emulating Starlink Gen2 and 3GPP-compliant measurement procedures, targets the critical challenge of selecting the optimal serving GW under rapidly time-varying LEO geometry to maximize network-wide utility.

To solve this high-dimensional optimization problem, this thesis developed a novel multi-agent dueling DDQN algorithm enhanced with PER and Polyak averaging. The ground GW agents learn to make cooperative decisions based on locally observable metrics—including satellite ephemeris, L3-filtered RSRP, and real-time traffic load. The reward function is designed to jointly optimize conflicting objectives, encoding a multi-objective trade-off that maximizes RSRP, ToS, and load fairness while minimizing switching instability (FLSO/UFLSO). Conceptually, this approach evolves the conventional link-centric handover triggering into a load-aware, learning-based policy that dynamically adapts to network conditions.

Extensive simulations confirm that the proposed policy achieves the most robust performance trade-off compared to both 3GPP-inspired heuristics and standard RL baselines. Unlike signal-based methods that suffer from excessive ping-pong effects or distance-based schemes that lead to load concentration and poor signal quality, the proposed agent successfully learns a balanced policy. It maintains high link quality and connection stability comparable to conservative heuristics, while simultaneously achieving superior load fairness and maximizing total traffic throughput. Specifically, the results demonstrate that the proposed method effectively unlocks the latent capacity of the ground segment by intelligently distributing traffic to underutilized resources without causing congestion.

Overall, this study verifies that a lightweight, decentralized RL-based controller can operate efficiently within realistic measurement periodicities. By realizing a favorable trade-off among switching stability, backhaul link quality, and network efficiency, the proposed framework offers a promising direction for autonomous and robust mobility management in future 6G NTN systems.